\section{Riemann-Hurwitz}
\label{section-riemann-hurewitz}

\noindent
Let $k$ be a field. Let $f : X \to Y$ be a morphism of smooth curves over $k$.
Then we obtain a canonical exact sequence
$$
f^*\Omega_{Y/k} \xrightarrow{\text{d}f} \Omega_{X/k}
\longrightarrow \Omega_{X/Y} \longrightarrow 0
$$
by Morphisms, Lemma \ref{morphisms-lemma-triangle-differentials}.
Since $X$ and $Y$ are smooth, the sheaves $\Omega_{X/k}$ and
$\Omega_{Y/k}$ are invertible modules, see
Morphisms, Lemma \ref{morphisms-lemma-smooth-omega-finite-locally-free}.
Assume the first map is nonzero, i.e., assume $f$ is generically
\'etale, see Lemma \ref{lemma-generically-etale}. Let $R \subset X$
be the closed subscheme cut out by the different $\mathfrak{D}_f$ of $f$.
By Discriminants, Lemma
\ref{discriminant-lemma-discriminant-quasi-finite-morphism-smooth}
this is the same as the vanishing locus of $\text{d}f$, it is
an effective Cartier divisor, and we get
$$
f^*\Omega_{Y/k} \otimes_{\mathcal{O}_X} \mathcal{O}_X(R) = \Omega_{X/k}
$$
In particular, if $X$, $Y$ are projective with
$k = H^0(Y, \mathcal{O}_Y) = H^0(X, \mathcal{O}_X)$
and $X$, $Y$ have genus $g_X$, $g_Y$, then we get the
Riemann-Hurwitz formula
\begin{align*}
2g_X - 2 & =
\deg(\Omega_{X/k}) \\
& =
\deg(f^*\Omega_{Y/k} \otimes_{\mathcal{O}_X} \mathcal{O}_X(R)) \\
& =
\deg(f) \deg(\Omega_{Y/k}) + \deg(R) \\
& =
\deg(f) (2g_Y - 2) + \deg(R)
\end{align*}
The first and last equality by Lemma \ref{lemma-genus-smooth}.
The second equality by the isomorphism of invertible sheaves given above.
The third equality by additivity of degrees
(Varieties, Lemma \ref{varieties-lemma-degree-tensor-product}),
the formula for the degree of a pullback
(Varieties, Lemma \ref{varieties-lemma-degree-pullback-map-proper-curves}),
and finally the formula for the degree of $\mathcal{O}_X(R)$
(Varieties, Lemma \ref{varieties-lemma-degree-effective-Cartier-divisor}).

\medskip\noindent
To use the Riemann-Hurwitz formula we need to compute
$\deg(R) = \dim_k \Gamma(R, \mathcal{O}_R)$. By the structure
of zero dimensional schemes over $k$ (see for example
Varieties, Lemma \ref{varieties-lemma-algebraic-scheme-dim-0}),
we see that $R$ is a finite disjoint union of spectra of
Artinian local rings $R = \coprod_{x \in R} \Spec(\mathcal{O}_{R, x})$
with each $\mathcal{O}_{R, x}$ of finite dimension over $k$. Thus
$$
\deg(R) = \sum\nolimits_{x \in R} \dim_k \mathcal{O}_{R, x} =
\sum\nolimits_{x \in R} d_x [\kappa(x) : k]
$$
with
$$
d_x = \text{length}_{\mathcal{O}_{R, x}} \mathcal{O}_{R, x} =
\text{length}_{\mathcal{O}_{X, x}} \mathcal{O}_{R, x}
$$
the multiplicity of $x$ in $R$
(see Algebra, Lemma \ref{algebra-lemma-pushdown-module}).
Let $x \in X$ be a closed point with image $y \in Y$.
Looking at stalks we obtain an exact sequence
$$
\Omega_{Y/k, y} \to \Omega_{X/k, x} \to \Omega_{X/Y, x} \to 0
$$
Choosing local generators $\eta_x$ and $\eta_y$ of the
(free rank $1$) modules $\Omega_{X/k, x}$ and $\Omega_{Y/k, y}$
we see that
$
\eta_y \mapsto h \eta_x
$
for some nonzero $h \in \mathcal{O}_{X, x}$. By the exact sequence we see that
$\Omega_{X/Y, x} \cong \mathcal{O}_{X, x}/h\mathcal{O}_{X, x}$
as $\mathcal{O}_{X, x}$-modules. Since the divisor
$R$ is cut out by $h$ (see above) we have
$\mathcal{O}_{R, x} = \mathcal{O}_{X, x}/h\mathcal{O}_{X, x}$.
Thus we find the following equalities
\begin{align*}
d_x
& =
\text{length}_{\mathcal{O}_{X, x}}(\mathcal{O}_{R, x}) \\
& =
\text{length}_{\mathcal{O}_{X, x}}(\mathcal{O}_{X, x}/h\mathcal{O}_{X, x}) \\
& =
\text{length}_{\mathcal{O}_{X, x}}(\Omega_{X/Y, x}) \\
& =
\text{ord}_{\mathcal{O}_{X, x}}(h) \\
& =
\text{ord}_{\mathcal{O}_{X, x}}(``\eta_y/\eta_x")
\end{align*}
The first equality by our definition of $d_x$. The second and third
we saw above. The fourth equality is the definition of $\text{ord}$, see
Algebra, Definition \ref{algebra-definition-ord}. Note that since
$\mathcal{O}_{X, x}$ is a discrete valuation ring, the integer
$\text{ord}_{\mathcal{O}_{X, x}}(h)$ just the valuation of $h$.
The fifth equality is a mnemonic.

\medskip\noindent
Here is a case where one can ``calculate'' the multiplicity $d_x$ in terms
of other invariants. Namely, if $\kappa(x)$ is separable over $k$, then
we may choose $\eta_x = \text{d}s$ and $\eta_y = \text{d}t$ where $s$
and $t$ are uniformizers in $\mathcal{O}_{X, x}$ and $\mathcal{O}_{Y, y}$
(Lemma \ref{lemma-uniformizer-works}).
Then $t \mapsto u s^{e_x}$ for some unit $u \in \mathcal{O}_{X, x}$
where $e_x$ is the ramification index of the extension
$\mathcal{O}_{Y, y} \subset \mathcal{O}_{X, x}$. Hence we get
$$
\eta_y = \text{d}t = \text{d}(u s^{e_x}) =
e s^{e_x - 1} u \text{d}s + s^{e_x} \text{d}u
$$
Writing $\text{d}u = w \text{d}s$ for some $w \in \mathcal{O}_{X, x}$
we see that
$$
``\eta_y/\eta_x" = e s^{e_x - 1} u + s^{e_x} w = (e_x u + s w)s^{e_x - 1}
$$
We conclude that the order of vanishing of this is $e_x - 1$
unless the characteristic of $\kappa(x)$ is $p > 0$ and $p$ divides $e_x$
in which case the order of vanishing is $> e_x - 1$.

\medskip\noindent
Combining all of the above we find that if $k$ has characteristic
zero, then
$$
2g_X - 2 = (2g_Y - 2)\deg(f) +
\sum\nolimits_{x \in X} (e_x - 1)[\kappa(x) : k]
$$
where $e_x$ is the ramification index of $\mathcal{O}_{X, x}$ over
$\mathcal{O}_{Y, f(x)}$. This precise formula will hold if and only
if all the ramification is tame, i.e., when the
residue field extensions $\kappa(x)/\kappa(y)$ are separable and
$e_x$ is prime to the characteristic of $k$, although the
arguments above are insufficient to prove this. We refer the reader
to Lemma \ref{lemma-rhe} and its proof.

